\appendix

\section[Basics on $q$-analysis]{Basics on $\boldsymbol{q}$-analysis}\label{app:qtheory}


A good reference for the material on $q$-analysis is \cite[Chapter~10]{AAR}. Assume $|q| < 1$ is an arbitrary complex number. Most statements work if $q$ is an indeterminate too.

The \textit{$q$-numbers} and the \textit{$q$-factorial} are def\/ined by
\begin{gather*}
[n]_q \myeq \frac{1 - q^n}{1 - q}, \qquad n\in\N,\\
[n]_q! \myeq [n]_q\cdots[2]_q[1]_q, \qquad n\in\N, \qquad [0]_q! \myeq 1.
\end{gather*}
It is evident that $[n]_q \rightarrow n$ and $[n]_q!\rightarrow n!$, as $q\rightarrow 1$, for any $n\in\Z_{\geq 0}$. Observe that we can also def\/ine $[x]_q$ for any $x\in\C$, as before, and it also holds that $[x]_q \rightarrow x$ as $q\rightarrow 1$. The \textit{$q$-Gamma function} is def\/ined by
\begin{gather*}
\Gamma_q(z) \myeq (1 - q)^{1-z}\frac{(q; q)_{\infty}}{(q^z; q)_{\infty}}.
\end{gather*}
From the def\/inition, the $q$-functional equation
\begin{gather*}
\Gamma_q(z + 1) = [z]_q\Gamma_q(z), \qquad z\notin\{\dots, -2, -1, 0\},
\end{gather*}
is evident. The $q$-Gamma function is a meromorphic function with simple poles at $z = 0, -1, -2, \dots$ and all their shifts by an integral multiple of $2\pi\sqrt{-1}/\ln{q}$.
The $q$-Gamma function has no zeroes in~$\C$. Moreover, we have the following convergence to the Gamma function.

\begin{thm}[{\cite[Corollary~10.3.4]{AAR}}]\label{qtheory1}
For any $z\in\C \setminus \{\dots, -2, -1, 0\}$, we have
\begin{gather*}
\lim_{q\rightarrow 1}{\Gamma_q(z)} = \Gamma(z).
\end{gather*}
\end{thm}
\begin{rem}\label{qtheory1remark}
As a consequence of Stieltjes--Vitali theorem, the convergence $\lim\limits_{q\rightarrow 1}{\Gamma_q(z)} = \Gamma(z)$ holds uniformly on compact subsets of $\C \setminus \{\dots, -2, -1, 0\}$.
\end{rem}

Other important identities we use in our paper are the \textit{$q$-binomial theorems}. To state them, we need to def\/ine the \textit{$q$-Pochhammer symbols} $(x; q)_n$ and $(x; q)_{\infty}$, for any $x\in\C$ and $n\in\Z_{\geq 0}$ by
\begin{gather*}
(x; q)_n \myeq
 \begin{cases}
 1 & \textrm{if } n = 0,\\
 \displaystyle \prod\limits_{i=1}^n{\big(1 - xq^{i-1}\big)} & \textrm{if } n\geq 1,\\
 \displaystyle \prod\limits_{i=1}^{\infty}{\big(1 - xq^{i-1}\big)} & \textrm{if } n = \infty.
 \end{cases}
\end{gather*}
Note that $|q| < 1$ implies that the product def\/ining $(x; q)_{\infty}$ is uniformly convergent for $x\in\C$, and thus $(x; q)_{\infty}$ is an entire function.

The $q$-binomial formula is the following

\begin{thm}[{\cite[Theorem~10.2.1]{AAR}}]\label{qtheory2}
For $|z| < 1$,
\begin{gather*}
\sum_{n=0}^{\infty}{\frac{(a; q)_n}{(q; q)_n}z^n} = \frac{(az; q)_{\infty}}{(z; q)_{\infty}}
\end{gather*}
\end{thm}

\begin{cor}\label{qtheory25}
For $z\in\C$, $m\in\N$,
\begin{gather*}
\sum_{n=0}^M{\frac{(q^{-1}; q^{-1})_M}{(q^{-1}; q^{-1})_n(q^{-1}; q^{-1})_{M-n}}(-1)^nq^{-{n \choose 2}}z^n} = \big(z; q^{-1}\big)_M.
\end{gather*}
\end{cor}
\begin{proof}
Let $a = q^{-M}$ in Theorem \ref{qtheory2}, and use $(q; q)_n = (-1)^n q^{{n + 1 \choose 2}} (q^{-1}; q^{-1})_n$. The statement is then proved for any $|z| < 1$; therefore it also holds for any $z\in\C$ because both sides are polynomials on $z$.
\end{proof}
